\section{多重共线性}

\begin{definition}[多重共线性]
    回归模型中，两个或多个自变量特征之间存在高度相关的关系。例如身高（厘米）和身高（英寸），这两个是完全的多重共线性。

    模型无法确定哪一个特征才是贡献率最大的。多重共线性会导致模型系数的方差很大，略微改变训练数据，各个相关特征的系数就会发生天翻地覆的变化。而且由于系数不稳定，无法再相信模型给出的系数大小，显著性检验也可能失去意义。
\end{definition}

\begin{theorem}[检测多重共线性]
    相关性矩阵计算所有自变量之间的相关系数，但只能检测两个变量之间的成对共线性，无法检测多变量复合的情况。
\end{theorem}

\begin{definition}[方差膨胀因子]

\end{definition}

\section{岭回归}

\begin{definition}[岭回归]
    岭回归是一种用于解决多重共线性问题的线性回归模型。当自变量之间存在高度相关性时，普通最小二乘法（OLS）估计会变得不稳定，导致模型系数的方差增大，预测精度下降。岭回归通过在损失函数中添加L2正则化项来解决这个问题，从而缩小系数的估计值，提高模型的泛化能力。
\end{definition}

\section{数学公式}
普通最小二乘法旨在最小化残差平方和：
$$ \min_{\beta} \|y - X\beta\|^2 $$
其中，$y$ 是响应变量，$X$ 是设计矩阵，$\beta$ 是回归系数向量。

岭回归的目标是最小化带有L2正则化项的目标函数：
$$ \min_{\beta} \left( \|y - X\beta\|^2 + \lambda \|\beta\|^2 \right) $$
其中，$\lambda \ge 0$ 是正则化参数（也称为收缩参数），用于控制正则化的强度。$\|\beta\|^2 = \sum_{j=1}^{p} \beta_j^2$ 是回归系数的L2范数。

当 $\lambda = 0$ 时，岭回归退化为普通最小二乘法。随着 $\lambda$ 的增大，回归系数 $\beta$ 的估计值会被压缩，趋向于零，但通常不会精确地等于零（与Lasso回归不同）。

\section{岭回归的优势}
\begin{itemize}
    \item \textbf{处理多重共线性}: 岭回归能够有效地处理自变量之间高度相关的情况，使模型更加稳定。
    \item \textbf{提高泛化能力}: 通过收缩系数，岭回归可以减少模型的方差，从而提高在未见过数据上的预测性能。
    \item \textbf{系数收缩}: 岭回归会使系数向零收缩，但通常不会使系数变为零，保留了所有变量的信息。
\end{itemize}

\section{岭回归的劣势}
\begin{itemize}
    \item \textbf{不进行特征选择}: 与Lasso回归不同，岭回归不会将系数精确地缩减到零，因此它不会进行特征选择。这意味着即使某些变量对预测不重要，它们仍然会保留在模型中。
    \item \textbf{正则化参数选择}: 需要通过交叉验证等方法来选择最优的正则化参数 $\lambda$。
\end{itemize}

\section{总结}
岭回归是一种强大的正则化技术，特别适用于存在多重共线性的回归问题。通过引入L2正则化，它能够稳定模型估计，提高预测精度，并增强模型的泛化能力。然而，它不具备特征选择的能力，这一点需要与Lasso回归等其他正则化方法进行权衡。

